Has broadband become a necessity good immune to price changes? Using a 15-year panel of 33 countries---27 European Union (EU) member states and 6 Eastern Partnership (EaP) countries (2010--2024)---and two-way fixed effects with Driscoll--Kraay standard errors, we document a fundamental transformation in broadband demand. Pre-COVID, EaP countries exhibited relatively elastic demand ($\varepsilon = \EaPBaseBr$, p$<$0.001)---a 10\% price reduction increased subscriptions by $\EaPTenPctChg$\%---while EU countries showed lower elasticity ($\varepsilon = \EUBaseBr$, p$<$0.10). By 2020--2024, both regions converged to near-zero elasticity, with price changes having no detectable effect on adoption. Crucially, placebo tests reveal this transformation began in 2015, not 2020, indicating a decade-long digital integration process rather than a COVID-19 shock. We further demonstrate that price measurement critically affects inference: income-relative prices (as \% of gross national income, GNI) yield significant EaP elasticities in all specifications, whereas purchasing-power-parity (PPP) adjusted prices yield significant results in only \PctPPPSigEaP\% of EaP specifications; for EU countries, all three price measures perform broadly similarly. These findings have immediate policy relevance: as broadband transitions from discretionary service to essential utility, policy emphasis must shift from affordability subsidies to universal infrastructure deployment.
